% c09 原书 p.104–116（PDF p.111–123）
% 接缝：p.103–104 属同一句，其上式与交换图已由 c08 收尾；本块从 p.104 的新段落开始。

对每个 $\cO_X$-模 $\cF$，定义 $\cF$ 关于 $\cI$ 的\emph{形式完全化} $\widehat\cF$ 为 $\cO_{\widehat X}$-模
\[
\widehat\cF=\plim_n(\cF/\cI^n\cF)|_{\widehat X}.
\]
我们有
\[
\widehat\cF=\plim_n\hat i_{n*}\hat i_n^*\cF,
\]
且典范态射族 $\cF\to i_{n*}i_n^*\cF$（$n\in\bN$）诱导态射
\[
\cF\to i_*\widehat\cF=\plim_n i_{n*}i_n^*\cF.
\]

设 $\cI'$ 是 $\cO_X$ 的凝聚理想层，使 $\Supp(\cO_X/\cI)=\Supp(\cO_X/\cI')$。则由 \ref{it:1.4.14}，存在自然数 $k$ 与 $k'$，使 $\cI^k\subset\cI'$ 且 $\cI'^{k'}\subset\cI$。所以 $X$ 与 $\cF$ 关于 $\cI$ 的形式完全化分别同构于 $X$ 与 $\cF$ 关于 $\cI'$ 的形式完全化。因此关于 $\cI$ 的形式完全化只依赖于以 $\cI$ 为理想层的 $X$ 的闭子概形的底拓扑空间。

\begin{proposition}\label{it:1.5.12}
设 $X$ 是诺特概形，$\cI$ 是 $\cO_X$ 的凝聚理想层，$(\widehat X,\cO_{\widehat X})$ 是 $X$ 关于 $\cI$ 的形式完全化。则 $(\widehat X,\cO_{\widehat X})$ 是局部环层空间，且典范态射 $(i,i^\#)\colon(\widehat X,\cO_{\widehat X})\to(X,\cO_X)$ 是局部环层空间的态射。
\end{proposition}

\begin{proof}
问题是局部的。可设 $X=\Spec A$，其中 $A$ 是诺特环。设 $I$ 是与 $\cI$ 对应的 $A$ 的理想。我们有
\[
\widehat X=\Supp(A/I)^{\sim}=V(I).
\]
$\widehat X$ 中的点对应于 $A$ 的满足 $\fp\supset I$ 的素理想 $\fp$。注意 $D(f)\cap\widehat X$（$f\in A-\fp$）构成 $\widehat X$ 中点 $\fp$ 的基本邻域系。所以
\[
\cO_{\widehat X,\fp}\cong\plim_{f\notin\fp}\cO_{\widehat X}\bigl(D(f)\cap\widehat X\bigr).
\]
我们有
\[
\cO_{\widehat X}\bigl(D(f)\cap\widehat X\bigr)\cong\plim_n(A/I^n)_f\cong\plim_n A_f/(I_f)^n\cong\widehat{A_f},
\]
其中 $\widehat{A_f}$ 是 $A_f$ 的 $I_f$-进完全化。所以
\[
\cO_{\widehat X,\fp}\cong\plim_{f\notin\fp}\widehat{A_f}.
\]
我们要证它是局部环。令
\[
\widehat{\fp_f}=\plim_n\fp_f/(I_f)^n.
\]
由 \ref{it:1.5.3}，它同构于 $\fp_f$ 的 $I_f$-进完全化。它是 $\widehat{A_f}$ 的理想。所以 $\plim_{f\notin\fp}\widehat{\fp_f}$ 是 $\plim_{f\notin\fp}\widehat{A_f}$ 的理想。要证 $\cO_{\widehat X,\fp}$ 是局部环，只需证明 $\plim_{f\notin\fp}\widehat{A_f}-\plim_{f\notin\fp}\widehat{\fp_f}$ 中每个元素都是单位。为此只需证明：对每个元素 $z\in\widehat{A_f}-\widehat{\fp_f}$，存在 $g\in A-\fp$，使得 $f$ 整除 $g$，且 $z$ 在典范同态 $\widehat{A_f}\to\widehat{A_g}$ 下的像是单位。设 $z_0$ 是 $z$ 在典范同态 $\widehat{A_f}\to A_f/I_f$ 下的像。选取 $\frac{x}{f^k}\in A_f$，使其在 $A_f/I_f$ 中的像为 $z_0$。由于 $z\notin\widehat{\fp_f}$，有 $x\notin\fp$。令 $g=fx$。则 $g\notin\fp$，且 $z_0$ 在典范同态 $A_f/I_f\to A_g/I_g$ 下的像 $z'_0$ 是单位。（其逆是 $A_g$ 中元素 $\frac{f^{k+1}}{g}$ 在 $A_g/I_g$ 中的像。）设 $z'$ 是 $z$ 在同态 $\widehat{A_f}\to\widehat{A_g}$ 下的像。则 $z'$ 在 $\widehat{A_g}\to A_g/I_g$ 下的像是 $z'_0$，而 $z'_0$ 是单位。由于 $\widehat{A_g}/\widehat{I_g}\cong A_g/I_g$，$z'$ 在 $\widehat{A_g}/\widehat{I_g}$ 中的像是单位，因而不属于 $\widehat{A_g}/\widehat{I_g}$ 的任何极大理想。但由 \ref{it:1.5.7}，$\widehat{I_g}\subset\operatorname{rad}(\widehat{A_g})$，所以 $z'$ 不属于 $\widehat{A_g}$ 的任何极大理想，从而是单位。这就证明了 $\plim_{f\notin\fp}\widehat{A_f}$ 是以 $\plim_{f\notin\fp}\widehat{\fp_f}$ 为极大理想的局部环。注意 $\plim_{f\notin\fp}A_f$ 的极大理想是 $\plim_{f\notin\fp}\fp_f$。容易看出典范同态
\[
\plim_{f\notin\fp}A_f\to\plim_{f\notin\fp}\widehat{A_f}
\]
是局部的。我们有 $\cO_{X,\fp}\cong\plim_{f\notin\fp}A_f$ 与 $\cO_{\widehat X,\fp}\cong\plim_{f\notin\fp}\widehat{A_f}$。所以典范态射 $(i,i^\#)\colon(\widehat X,\cO_{\widehat X})\to(X,\cO_X)$ 是局部环层空间的态射。
\end{proof}
\begin{tnote}
证明的关键是把完全化与局部化交换：先在仿射开集上作 $I$-进完全化，再取 $\fp$ 处的茎。$\widehat X$ 上拓扑空间的点是 $V(I)$ 中的素理想，而结构层的茎则是 $\widehat{A_f}$ 沿 $f\notin\fp$ 的逆向极限。局部性最终来自 \ref{it:1.5.7} 的 $\widehat I\subset\operatorname{rad}(\widehat A)$。[依据：原书 p.105–106 命题 1.5.12 的证明]
\end{tnote}

\begin{proposition}\label{it:1.5.13}
设 $X$ 是诺特概形，$\cI$ 是 $\cO_X$ 的凝聚理想层，$\widehat X$ 是 $X$ 关于 $\cI$ 的形式完全化。
\begin{enumerate}[label=(\roman*),leftmargin=2.8em]
  \item 对每个凝聚 $\cO_X$-模的正合序列
  \[
  0\to\cF'\to\cF\to\cF''\to 0,
  \]
  序列
  \[
  0\to\widehat{\cF'}\to\widehat\cF\to\widehat{\cF''}\to 0
  \]
  正合。
  \item 设 $i\colon\widehat X\to X$ 是典范态射。则对每个凝聚 $\cO_X$-模 $\cF$，有典范同构
  \[
  i^*\cF\cong\widehat\cF.
  \]
  \item 对凝聚 $\cO_X$-模 $\cF$ 与 $\cG$，有
  \[
  \widehat\cF\tensor_{\cO_{\widehat X}}\widehat\cG\cong\widehat{(\cF\tensor_{\cO_X}\cG)},\qquad
  \Hom_{\cO_{\widehat X}}(\widehat\cF,\widehat\cG)\cong\widehat{(\Hom_{\cO_X}(\cF,\cG))}.
  \]
\end{enumerate}
\end{proposition}

\begin{proof}
(i) 设
\[
0\to\cF'\to\cF\to\cF''\to 0
\]
是凝聚 $\cO_X$-模的正合序列。设 $U=\Spec A$ 是 $X$ 的仿射开子概形。设对某个有限生成 $A$-模 $M'$、$M$ 与 $M''$，有 $\cF'|_U\cong M'^{\sim}$，$\cF|_U\cong M^{\sim}$，$\cF''|_U\cong M''^{\sim}$。则有正合序列
\[
0\to M'\to M\to M''\to 0.
\]
设 $I$ 是与 $\cI$ 对应的 $A$ 的理想。我们有
\[
\widehat\cF(U\cap\widehat X)=\plim_n(\cF/\cI^n\cF)|_{\widehat X}(U\cap\widehat X)\cong\plim_n M/I^nM=\widehat M,
\]
其中 $\widehat M$ 是 $M$ 的 $I$-进完全化。类似地，有 $\widehat\cF'(U\cap\widehat X)\cong\widehat{M'}$ 与 $\widehat\cF''(U\cap\widehat X)\cong\widehat{M''}$。由 \ref{it:1.5.4}，序列
\[
0\to\widehat\cF'(U\cap\widehat X)\to\widehat\cF(U\cap\widehat X)\to\widehat\cF''(U\cap\widehat X)\to 0
\]
正合。当 $U$ 取遍 $X$ 的仿射开子概形族时，$U\cap\widehat X$ 构成 $\widehat X$ 的拓扑的一组基。所以序列
\[
0\to\widehat{\cF'}\to\widehat\cF\to\widehat{\cF''}\to 0
\]
正合。

(ii) 典范态射 $\cF\to i_*\widehat\cF$ 诱导态射 $i^*\cF\to\widehat\cF$。下面证明它是同构。这是局部问题，故可设存在正合序列
\[
\cO_X^m\to\cO_X^n\to\cF\to 0.
\]
考虑交换图
\begin{center}
\begin{tikzcd}
i^*\cO_X^m \arrow[r] \arrow[d] & i^*\cO_X^n \arrow[r] \arrow[d] & i^*\cF \arrow[r] \arrow[d] & 0\\
\widehat{\cO_X}^{\,m} \arrow[r] & \widehat{\cO_X}^{\,n} \arrow[r] & \widehat\cF \arrow[r] & 0.
\end{tikzcd}
\end{center}
两行都正合。前两个竖直箭头是同构。所以第三个竖直箭头也是同构。

(iii) 第一个等式由 (ii) 及以下事实推出：
\[
i^*\cF\tensor_{i^{-1}\cO_X}i^*\cG\cong i^*(\cF\tensor_{\cO_X}\cG).
\]
有典范态射
\[
\Hom_{\cO_X}(\cF,\cG)\to i_*\Hom_{\cO_{\widehat X}}(i^*\cF,i^*\cG).
\]
它诱导态射
\[
i^*\Hom_{\cO_X}(\cF,\cG)\to\Hom_{\cO_{\widehat X}}(i^*\cF,i^*\cG).
\]
下面证明它是同构。可局部地做，并设存在正合序列
\[
\cO_X^m\to\cO_X^n\to\cF\to 0.
\]
考虑交换图
\begin{center}
\begin{tikzcd}
0 \arrow[r] & i^*\Hom_{\cO_X}(\cF,\cG) \arrow[r] \arrow[d] & i^*\Hom_{\cO_X}(\cO_X^n,\cG) \arrow[r] \arrow[d] & i^*\Hom_{\cO_X}(\cO_X^m,\cG) \arrow[d]\\
0 \arrow[r] & \Hom_{\cO_{\widehat X}}(i^*\cF,i^*\cG) \arrow[r] & \Hom_{\cO_{\widehat X}}(i^*\cO_X^n,i^*\cG) \arrow[r] & \Hom_{\cO_{\widehat X}}(i^*\cO_X^m,i^*\cG).
\end{tikzcd}
\end{center}
由 (i) 与 (ii)，两行都正合，且最后两个竖直箭头是同构。所以第一个竖直箭头是同构。我们的断言于是由 (ii) 推出。
\end{proof}
\begin{tnote}
\ref{it:1.5.13} 说明：对凝聚模而言，形式完全化函子 $\cF\mapsto\widehat\cF$ 正合，并且保持张量积与 $\Hom$。这与代数侧的 \ref{it:1.5.4} 互相呼应。注意 (i) 中对一般（非凝聚）层并不成立。[依据：原书 p.106–107 命题 1.5.13 及其证明]
\end{tnote}

\begin{proposition}\label{it:1.5.14}
设 $X$ 是诺特概形，$\cI$ 是 $\cO_X$ 的凝聚理想层，$\widehat X$ 是 $X$ 关于 $\cI$ 的形式完全化。
\begin{enumerate}[label=(\roman*),leftmargin=2.8em]
  \item 对 $\cO_X$ 的每个凝聚理想层 $\cJ$，$\cJ$ 关于 $\cI$ 的形式完全化 $\widehat\cJ$ 是 $\cO_{\widehat X}$ 的理想层。
  \item 对每个 $n$，有 $\cO_{\widehat X}/\widehat\cI^n\cong(\cO_X/\cI^n)|_{\widehat X}$。
  \item 对每对 $0\le m\le n$，有 $\widehat\cI^n\subset\widehat\cI^m$ 且 $\widehat\cI^m/\widehat\cI^n\cong(\cI^m/\cI^n)|_{\widehat X}$。
  \item 对 $\cO_X$ 的每个凝聚理想层 $\cJ$ 与每个 $n$，有
  \[
  \widehat{\cJ^n}\cong i^*(\cJ^n)\cong\cJ^n\cO_{\widehat X}\cong\widehat\cJ^{\,n}.
  \]
\end{enumerate}
\end{proposition}

\begin{proof}
有正合序列
\[
0\to\cJ\to\cO_X.
\]
由 \ref{it:1.5.13} (i)，序列
\[
0\to\widehat\cJ\to\widehat{\cO_X}
\]
正合，即
\[
0\to\widehat\cJ\to\cO_{\widehat X}
\]
正合。所以 $\widehat\cJ$ 是 $\cO_{\widehat X}$ 的理想层。

对每对 $0\le m\le n$，有凝聚 $\cO_X$-模的正合序列
\[
0\to\cI^n\to\cI^m\to\cI^m/\cI^n\to 0.
\]
由 \ref{it:1.5.13} (i)，序列
\[
0\to\widehat{\cI^n}\to\widehat{\cI^m}\to\widehat{\cI^m/\cI^n}\to 0
\]
正合。由形式完全化的定义可证
\[
\widehat{\cI^m/\cI^n}\cong(\cI^m/\cI^n)|_{\widehat X}.
\]
这说明 $\widehat\cI^n\subset\widehat\cI^m$ 且
\[
\widehat\cI^m/\widehat\cI^n\cong(\cI^m/\cI^n)|_{\widehat X}.
\]
取 $m=0$，得
\[
\cO_{\widehat X}/\widehat\cI^n\cong(\cO_X/\cI^n)|_{\widehat X}.
\]
由 \ref{it:1.5.13} (i) 与 (ii)，典范态射
\[
i^*\cJ^n\to i^*\cO_X\cong\cO_{\widehat X}
\]
是单射。其像是 $\cJ^n\cO_{\widehat X}$。所以
\[
i^*\cJ^n\cong\cJ^n\cO_{\widehat X}.
\]
(iv) 中的其余断言于是显然。
\end{proof}

\begin{proposition}\label{it:1.5.15}
设 $X$ 是诺特概形，$\cI$ 是 $\cO_X$ 的凝聚理想层，$\cF$ 是凝聚 $\cO_X$-模，$\widehat X$ 与 $\widehat\cF$ 分别是 $X$ 与 $\cF$ 关于 $\cI$ 的形式完全化。则典范同态 $\cF(X)\to\widehat\cF(\widehat X)$ 的核，由那些在 $X$ 中 $\widehat X$ 的某个邻域上消失的截面 $s\in\cF(X)$ 组成。
\end{proposition}

\begin{proof}
若 $s\in\cF(X)$ 在 $X$ 中 $\widehat X$ 的某个邻域上消失，则它在每个 $(\cF/\cI^n\cF)|_{\widehat X}(\widehat X)$ 中的像为 $0$。因此它在 $\widehat\cF(\widehat X)$ 中的像为 $0$。反之，设 $s\in\ker(\cF(X)\to\widehat\cF(\widehat X))$。则对 $X$ 的每个仿射开子概形 $U=\Spec A$，$s|_U$ 落在 $\cF(U)\to\widehat\cF(U\cap\widehat X)$ 的核中。设对某个有限生成 $A$-模 $M$ 有 $\cF|_U\cong M^{\sim}$，并设 $I$ 是与 $\cI$ 对应的 $A$ 的理想。我们有 $\cF(U)\cong M$ 与 $\widehat\cF(U\cap\widehat X)=\widehat M$，其中 $\widehat M$ 是 $M$ 的 $I$-进完全化。设 $x\in M$ 是对应于 $s|_U\in\cF(U)$ 的元素。则 $x$ 落在典范同态 $M\to\widehat M$ 的核中。由 \ref{it:1.5.8}，存在 $r\in I$，使 $(1+r)x=0$。注意 $1+r$ 在 $A_{\fp}$ 中的像对每个 $\fp\in V(I)\cong U\cap\widehat X$ 都是单位。所以存在 $V(I)$ 在 $U=\Spec A$ 中的邻域 $V$，使 $x$ 在 $M^{\sim}(V)$ 中的像为 $0$。由于 $U$ 是任意的，存在 $\widehat X$ 在 $X$ 中的邻域，使 $s$ 在其上消失。
\end{proof}

\begin{corollary}\label{it:1.5.16}
设 $X$ 是诺特概形，$\cI$ 是 $\cO_X$ 的凝聚理想层，$\widehat X$ 是 $X$ 关于 $\cI$ 的形式完全化，$u\colon\cF\to\cG$ 是凝聚 $\cO_X$-模的态射，$\hat u\colon\widehat\cF\to\widehat\cG$ 是 $u$ 在形式完全化上诱导的态射。
\begin{enumerate}[label=(\roman*),leftmargin=2.8em]
  \item $\hat u$ 消失，当且仅当 $u$ 在 $X$ 中 $\widehat X$ 的某个邻域上消失。
  \item $\hat u$ 是单态射（resp.\ 满态射，resp.\ 同构），当且仅当 $u$ 在 $X$ 中 $\widehat X$ 的某个邻域上是单态射（resp.\ 满态射，resp.\ 同构）。
\end{enumerate}
\end{corollary}

\begin{proof}
(i) 由 \ref{it:1.4.6}，$\Hom_{\cO_X}(\cF,\cG)$ 是凝聚的。把 $u$ 视为 $\Hom_{\cO_X}(\cF,\cG)$ 的截面，我们的断言由 \ref{it:1.5.15} 应用于此凝聚层推出。

(ii) 有正合序列
\[
0\to\ker u\to\cF\xrightarrow{\ u\ }\cG\to\coker u\to 0.
\]
由 \ref{it:1.5.13} (i)，序列
\[
0\to\widehat{\ker u}\to\widehat\cF\xrightarrow{\ \hat u\ }\widehat\cG\to\widehat{\coker u}\to 0
\]
正合。所以 $\hat u$ 是单态射，当且仅当 $\widehat{\ker u}\to\widehat\cF$ 消失。由 (i)，这等价于说 $\ker u\to\cF$ 在 $\widehat X$ 的某个邻域上消失，即 $u$ 在 $X$ 中 $\widehat X$ 的某个邻域上是单态射。类似地可证 (ii) 中的其余断言。
\end{proof}

设 $X$ 是诺特概形，$\cI$ 是 $\cO_X$ 的凝聚理想层，$\widehat X$ 是 $X$ 关于 $\cI$ 的形式完全化，$X_n$ 是由理想层 $\cI^n$（$n\in\bN$）定义的 $X$ 的闭子概形。用下面的交换图（$m\ge n$）固定记号。
\begin{center}
\begin{tikzcd}
(X_n,\cO_{X_n}) \arrow[r,"i_n"] \arrow[d,"i_{nm}"'] & (X,\cO_X)\\
(X_m,\cO_{X_m}) \arrow[r,"\hat i_m"'] \arrow[ur,"i_m"] & (\widehat X,\cO_{\widehat X}) \arrow[u,"i"']
\end{tikzcd}
\end{center}
其中所有箭头都是典范态射。\emph{凝聚逆向系统}\idx{ningjunixiangxitong}{凝聚逆向系统} $(\cF_n,\phi_{mn})_{n\in\bN}$ 由以下数据组成：对每个 $n$，$X_n$ 上的一个 $\cO_{X_n}$-模 $\cF_n$；对每对 $m\ge n$，一个态射 $\phi_{mn}\colon\cF_m\to i_{nm*}\cF_n$；并且以下条件成立：
\begin{enumerate}[label=(\alph*),leftmargin=2.8em]
  \item 每个 $\cF_n$ 是凝聚 $\cO_{X_n}$-模。
  \item 对每个 $n$ 有 $\phi_{nn}=\id_{\cF_n}$；且对每个三元组 $k\ge m\ge n$ 有 $i_{mk*}(\phi_{mn})\circ\phi_{km}=\phi_{kn}$。
  \item 对每对 $m\ge n$，$\phi_{mn}$ 诱导同构 $i_{nm}^*\cF_m\cong\cF_n$，即 $\cF_m\tensor_{\cO_{X_m}}\cO_{X_n}\cong\cF_n$。
\end{enumerate}

$\cO_{\widehat X}$-模 $\fF$ 称为\emph{凝聚的}，若对某个凝聚逆向系统 $(\cF_n,\phi_{mn})_{n\in\bN}$ 有
\[
\fF\cong\plim_n\hat i_{n*}\cF_n.
\]

设 $\cF$ 是凝聚 $\cO_X$-模。则 $(\hat i_n^*\cF,\phi_{mn})_{n\in\bN}$ 构成凝聚逆向系统，其中对每对 $m\ge n$，$\phi_{mn}$ 是典范态射
\[
\hat i_m^*\cF\to i_{nm*}\hat i_{nm}^*\hat i_m^*\cF\cong i_{nm*}\hat i_n^*\cF.
\]
所以
\[
\widehat\cF=\plim_n\hat i_{n*}\hat i_n^*\cF
\]
是凝聚 $\cO_{\widehat X}$-模。凝聚 $\cO_{\widehat X}$-模若同构于某个凝聚 $\cO_X$-模 $\cF$ 的 $\widehat\cF$，则称为\emph{可代数化}\idx{kedaishuhua}{可代数化}的。并非每个凝聚 $\cO_{\widehat X}$-模都是可代数化的。但 \ref{it:1.5.20} 与 \ref{it:2.7.10} 给出了答案为肯定的若干情形。
\begin{tnote}
「可代数化」是本节最易误解的概念之一：一个凝聚 $\cO_{\widehat X}$-模未必来自 $X$ 上的凝聚模，因为它只在 $\widehat X$ 上有定义，而 $\widehat X$ 之外的层信息已全部丢失。\ref{it:1.5.20} 说明当 $A$ 是 $I$-进完备的仿射概形时，情形总是好的；\ref{it:2.7.10} 则给出一般的充分条件，其证明正是本节 \ref{it:1.5.10}--\ref{it:1.5.21} 这些代数结论的用途。[依据：原书 p.110 的陈述及本书 §2.7 对此处的引用]
\end{tnote}

\begin{proposition}\label{it:1.5.17}
记号如上。设 $(\cF_n,\phi_{mn})_{n\in\bN}$ 是凝聚逆向系统，$\fF=\plim_n\hat i_{n*}\cF_n$。则对每个 $n$，有典范同构
\[
\hat i_n^*\fF=\fF\tensor_{\cO_{\widehat X}}\cO_{X_n}\cong\fF/\widehat\cI^n\fF\cong\cF_n,
\]
且对每对 $m\ge n$，$\phi_{mn}\colon\cF_m\to i_{nm*}\cF_n$ 可等同于典范态射
\[
\hat i_m^*\fF\to i_{nm*}\hat i_{nm}^*\hat i_m^*\fF\cong i_{nm*}\hat i_n^*\fF.
\]
\end{proposition}

\begin{proof}
每个典范态射 $\plim_n\hat i_{n*}\cF_n\to\hat i_{n*}\cF_n$ 诱导态射
\[
\hat i_n^*\fF\to\cF_n.
\]
下面证明它是同构。只需证明：对 $X$ 的每个仿射开子概形 $U=\Spec A$，典范同态
\[
\fF(U\cap\widehat X)\tensor_{\cO_{\widehat X}(U\cap\widehat X)}\cO_{X_n}(U\cap X_n)\to\cF_n(U\cap X_n)
\]
是同构。对每个 $n$，设 $M_n$ 是有限生成 $A/I^n$-模，使 $\cF_n|_{U\cap X_n}\cong M_n^{\sim}$。对每对 $m\ge n$，态射 $\phi_{mn}\colon\cF_m\to i_{nm*}\cF_n$ 诱导同态 $\psi_{mn}\colon M_m\to M_n$，它与模结构及环同态 $A/I^m\to A/I^n$ 相容。由于 $\phi_{mn}$ 诱导同构 $i_{nm}^*\cF_m\cong\cF_n$，$\psi_{mn}$ 是满射，其核为 $I^nM_m$。令 $M=\plim_n M_n$。由 \ref{it:1.5.11}，典范同态 $M\to M_n$ 诱导同构
\[
M\tensor_{\widehat A}A/I^n\cong M_n.
\]
但我们有
\begin{align*}
\fF(U\cap\widehat X)&=\plim_n\cF_n(U\cap X_n)\cong\plim_n M_n,\\
\cO_{\widehat X}(U\cap\widehat X)&\cong\widehat A,\\
\cO_{X_n}(U\cap X_n)&\cong A/I^n,\\
\cF_n(U\cap X_n)&\cong M_n,
\end{align*}
所以
\[
\fF(U\cap\widehat X)\tensor_{\cO_{\widehat X}(U\cap\widehat X)}\cO_{X_n}(U\cap X_n)\cong\cF_n(U\cap X_n).
\]
由 \ref{it:1.5.14}，有
\[
\hat i_n^*\fF=\fF\tensor_{\cO_{\widehat X}}\cO_{X_n}\cong\fF\tensor_{\cO_{\widehat X}}\cO_{\widehat X}/\widehat\cI^n\cO_{\widehat X}\cong\fF/\widehat\cI^n\fF.
\]
对每对 $m\ge n$，态射 $i_{nm}\colon X_n\to X_m$ 在底拓扑空间上是恒等映射。所以 $\phi_{mn}\colon\cF_m\to i_{nm*}\cF_n$ 可视为态射 $\cF_m\to\cF_n$。类似地，典范态射 $\hat i_m^*\fF\to i_{nm*}\hat i_{nm}^*\hat i_m^*\fF\cong i_{nm*}\hat i_n^*\fF$ 可视为态射 $\fF/\widehat\cI^m\fF\to\fF/\widehat\cI^n\fF$。有交换图
\begin{center}
\begin{tikzcd}
\fF/\widehat\cI^m\fF \arrow[r] \arrow[d] & \fF/\widehat\cI^n\fF \arrow[d]\\
\cF_m \arrow[r,"\phi_{mn}"'] & \cF_n.
\end{tikzcd}
\end{center}
所以 $\phi_{mn}$ 可等同于典范态射 $\hat i_m^*\fF\to i_{nm*}\hat i_n^*\fF$。
\end{proof}

\begin{proposition}\label{it:1.5.18}
保留上述记号。设 $\fF$ 是 $\cO_{\widehat X}$-模，$(\cF_n,\phi_{mn})_{n\in\bN}$ 是由 $\cF_n=\hat i_n^*\fF$（对每个 $n$）及以下典范态射（对每对 $m\ge n$）定义的逆向系统：
\[
\phi_{mn}\colon\hat i_m^*\fF\to i_{nm*}\hat i_{nm}^*\hat i_m^*\fF\cong i_{nm*}\hat i_n^*\fF.
\]
则 $\fF$ 是凝聚 $\cO_{\widehat X}$-模，当且仅当每个 $\cF_n$ 是凝聚 $\cO_{X_n}$-模且 $\fF\cong\plim_n\hat i_{n*}\cF_n$。
\end{proposition}

\begin{proof}
显然有 $i_{nm}^*\cF_m\cong\cF_n$。「当」部分由凝聚 $\cO_{\widehat X}$-模的定义推出。「仅当」部分由 \ref{it:1.5.17} 推出。
\end{proof}

\begin{corollary}\label{it:1.5.19}
记号如上。设 $(\cF_n,\phi_{mn})_{n\in\bN}$ 与 $(\cG_n,\psi_{mn})_{n\in\bN}$ 是两个凝聚逆向系统，$\fF=\plim_n\hat i_{n*}\cF_n$，$\fG=\plim_n\hat i_{n*}\cG_n$。则有双射
\[
\Hom_{\cO_{\widehat X}}(\fF,\fG)\cong\plim_n\Hom_{\cO_{X_n}}(\cF_n,\cG_n).
\]
\end{corollary}

\begin{proof}
$\plim_n\Hom_{\cO_{X_n}}(\cF_n,\cG_n)$ 中的每个元素 $(u_n)$ 定义一族态射
\[
\hat i_{n*}(u_n)\colon\hat i_{n*}\cF_n\to\hat i_{n*}\cG_n.
\]
取逆向极限，得到态射
\[
\plim_n\hat i_{n*}(u)\colon\plim_n\hat i_{n*}\cF_n\to\plim_n\hat i_{n*}\cG_n.
\]
定义 $S$ 为
\[
S\colon\plim_n\Hom_{\cO_{X_n}}(\cF_n,\cG_n)\to\Hom_{\cO_{\widehat X}}(\fF,\fG),\quad
(u_n)\mapsto\plim_n\hat i_{n*}(u).
\]
任意 $u\in\Hom_{\cO_{\widehat X}}(\fF,\fG)$ 诱导一族态射 $\hat i_n^*(u)\colon\hat i_n^*\fF\to\hat i_n^*\fG$。由 \ref{it:1.5.17}，这族态射可等同于一族态射 $u_n\colon\cF_n\to\cG_n$。容易验证 $(u_n)\in\plim_n\Hom_{\cO_{X_n}}(\cF_n,\cG_n)$。定义
\[
T\colon\Hom_{\cO_{\widehat X}}(\fF,\fG)\to\plim_n\Hom_{\cO_{X_n}}(\cF_n,\cG_n),\quad u\mapsto(u_n).
\]
容易验证 $S$ 与 $T$ 互为逆。
\end{proof}

\begin{proposition}\label{it:1.5.20}
设 $A$ 是诺特环，$I$ 是 $A$ 的理想，$X=\Spec A$，$\widehat X$ 是 $X$ 关于 $\cI^{\sim}$ 的形式完全化。若 $A$ 是 $I$-进完备的，则每个凝聚 $\cO_{\widehat X}$-模都形如 $(M^{\sim})^{\wedge}$，其中 $M$ 是某个有限生成 $A$-模，$(M^{\sim})^{\wedge}$ 是 $M^{\sim}$ 关于 $\cI^{\sim}$ 的形式完全化。
\end{proposition}

\begin{proof}
设 $\fF$ 是凝聚 $\cO_{\widehat X}$-模。则对某个凝聚逆向系统 $(\cF_n,\phi_{mn})_{n\in\bN}$ 有 $\fF\cong\plim_n\hat i_{n*}\cF_n$。在 $X_n=\Spec A/I^n$ 上，有 $\cF_n\cong M_n^{\sim}$，其中 $M_n$ 是某个有限生成 $A/I^n$-模。对每对 $m\ge n$，态射 $\phi_{mn}\colon\cF_m\to i_{nm*}\cF_n$ 诱导同态 $\psi_{mn}\colon M_m\to M_n$，它与模结构及环同态 $A/I^m\to A/I^n$ 相容。此外，由于 $\phi_{mn}$ 诱导同构 $i_{nm}^*\cF_m\cong\cF_n$，$\psi_{mn}$ 是满射，其核为 $I^nM_m$。令 $M=\plim_n M_n$。由 \ref{it:1.5.11}，$M$ 是 $\widehat A=\plim_n A/I^n$ 上的有限生成模，且典范同态 $M\to M_n$ 诱导同构
\[
M\tensor_{\widehat A}A/I^n\cong M_n.
\]
由于 $A$ 是 $I$-进完备的，$M$ 是有限生成 $A$-模，且典范同态 $M\to M_n$ 诱导同构
\[
M/I^nM\cong M_n.
\]
我们有
\begin{align*}
(M^{\sim})^{\wedge}
&=\plim_n(M^{\sim}/(\cI^{\sim})^nM^{\sim})|_{\widehat X}\\
&\cong\plim_n(M/I^nM)^{\sim}|_{\widehat X}\\
&\cong\plim_n\hat i_{n*}M_n^{\sim}\\
&\cong\plim_n\hat i_{n*}\cF_n\\
&\cong\fF,
\end{align*}
即 $\fF\cong(M^{\sim})^{\wedge}$。
\end{proof}

\begin{proposition}\label{it:1.5.21}
设 $X$ 是诺特概形，$\cI$ 是 $\cO_X$ 的凝聚理想层，$\widehat X$ 是 $X$ 关于 $\cI$ 的形式完全化，$\fF$ 与 $\fG$ 是凝聚 $\cO_{\widehat X}$-模。
\begin{enumerate}[label=(\roman*),leftmargin=2.8em]
  \item $\fF\tensor_{\cO_{\widehat X}}\fG$ 与 $\Hom_{\cO_{\widehat X}}(\fF,\fG)$ 是凝聚 $\cO_{\widehat X}$-模。
  \item 对每个态射 $u\colon\fF\to\fG$，$\cO_{\widehat X}$-模 $\ker u$、$\coker u$ 与 $\Image u$ 都是凝聚的。
\end{enumerate}
\end{proposition}

\begin{proof}
我们只证 (ii)，(i) 的证明类似。设 $\fK=\ker\phi$，并对每个 $n$ 令 $\fK_n=\hat i_n^*\fK$。
\begin{tnote}
原书 p.113 此处作「Let $\fK=\ker\phi$」，据上下文应为 $\fK=\ker u$，其中 $u\colon\fF\to\fG$ 是 (ii) 中给定的态射。[依据：原书 p.113 页图像与 1.5.21 (ii) 的陈述]
\end{tnote}
由 \ref{it:1.5.18}，需证每个 $\fK_n$ 是凝聚 $\cO_{X_n}$-模，且 $\fK\cong\plim_n\hat i_{n*}\fK_n$。这是局部问题。可设 $X=\Spec A$，其中 $A$ 是诺特环。设 $I$ 是与 $\cI$ 对应的 $A$ 的理想。注意把 $A$ 与 $I$ 换为它们的 $I$-进完全化 $\widehat A$ 与 $\widehat I$，不改变 $\widehat X$ 与 $X_n$，从而对 $\fK$ 与 $\fK_n$ 没有影响。作此替换后，可化归为 $A$ 是 $I$-进完备的情形。由 \ref{it:1.5.20}，有 $\fF\cong(M^{\sim})^{\wedge}$ 与 $\fG\cong(N^{\sim})^{\wedge}$，其中 $M$ 与 $N$ 是有限生成 $A$-模。由 \ref{it:1.5.19}，态射 $u\colon\fF\to\fG$ 由一族相容的同态 $u_n\colon M/I^nM\to N/I^nN$ 诱导。取逆向极限，并利用 $M$ 与 $N$ 由 \ref{it:1.5.5} 必为 $I$-进完备的，得到同态 $\hat u=\plim_n u_n\colon M\to N$。由 \ref{it:1.4.2} (iii) 与 \ref{it:1.5.13} (i)，序列
\[
0\to((\ker\hat u)^{\sim})^{\wedge}\to(M^{\sim})^{\wedge}\to(N^{\sim})^{\wedge}
\]
正合。所以 $\fK\cong((\ker\hat u)^{\sim})^{\wedge}$，从而 $\fK$ 是凝聚的。其余断言同理可证。
\end{proof}

\begin{proposition}\label{it:1.5.22}
设 $f\colon X\to Y$ 是诺特概形的态射，$\cI$ 是 $\cO_Y$ 的凝聚理想层，$\widehat X$ 与 $\widehat Y$ 分别是 $X$ 与 $Y$ 关于 $\cI\cO_X$ 与 $\cI$ 的形式完全化，$\hat f\colon\widehat X\to\widehat Y$ 是由 $f$ 诱导的局部环层空间态射。
\begin{enumerate}[label=(\roman*),leftmargin=2.8em]
  \item 对每个凝聚 $\cO_Y$-模 $\cG$，有典范同构
  \[
  \widehat{f^*\cG}\cong\hat f^*\widehat\cG,
  \]
  其中 $\widehat{f^*\cG}$ 与 $\widehat\cG$ 分别是 $f^*\cG$ 与 $\cG$ 关于 $\cI\cO_X$ 与 $\cI$ 的形式完全化。
  \item 对每个凝聚 $\cO_{\widehat Y}$-模 $\fG$，$\hat f^*\fG$ 是凝聚 $\cO_{\widehat X}$-模。
\end{enumerate}
\end{proposition}

\begin{proof}
(i) 有交换图
\begin{center}
\begin{tikzcd}
\widehat X \arrow[r,"\hat i_X"] \arrow[d,"\hat f"'] & X \arrow[d,"f"]\\
\widehat Y \arrow[r,"\hat i_Y"'] & Y,
\end{tikzcd}
\end{center}
其中水平箭头是典范态射。$\cG$ 与 $f^*\cG$ 都是凝聚的。由 \ref{it:1.5.13} (ii)，有 $\widehat\cG\cong\hat i_Y^*\cG$ 与 $\widehat{f^*\cG}\cong\hat i_X^*f^*\cG$。我们的断言于是由 $\hat i_X^*f^*\cG\cong\hat f^*\hat i_Y^*\cG$ 推出。

(ii) 由 \ref{it:1.5.18}，证明 $\hat f^*\fG$ 是凝聚的是局部问题。所以可设 $X=\Spec B$，$Y=\Spec A$，其中 $A$ 是诺特环，$B$ 是诺特 $A$-代数。设 $I$ 是与 $\cI$ 对应的 $A$ 的理想。把 $A$、$I$ 与 $B$ 换为它们的 $I$-进完全化 $\widehat A$、$\widehat I$ 与 $\widehat B$，不改变 $\widehat X$、$\widehat Y$ 与 $\hat f$。所以可设 $A$ 是 $I$-进完备的。由 \ref{it:1.5.20}，有 $\fG\cong(M^{\sim})^{\wedge}$，其中 $M$ 是有限生成 $A$-模。由 (i)，有
\[
\hat f^*\bigl((M^{\sim})^{\wedge}\bigr)\cong(f^*M^{\sim})^{\wedge}.
\]
所以 $\hat f^*\fG\cong(f^*M^{\sim})^{\wedge}$ 是凝聚的。
\end{proof}

\begin{exercises}
\item 设 $A$ 是环，$I$ 是 $A$ 的理想。$A$ 上的 \emph{$I$-进拓扑}是指使每个 $a\in A$ 的基本邻域系由 $a+I^n$（$n\in\bN$）组成的拓扑。$A$ 中的序列 $\{a_n\}$ 称为 \emph{Cauchy 序列}，若
\[
\lim_{m,n\to\infty}(a_n-a_m)=0.
\]
证明：Cauchy 序列构成环；以 $0$ 为极限的收敛序列构成此环的理想；此环对此理想的商同构于 $\widehat A=\plim_n A/I^n$。\label{ex:1.5.1}

\item 设 $A$ 是一个理想，$I$ 是 $A$ 的理想，$\widehat A$ 是 $A$ 关于 $I$-进拓扑的完全化。证明以下断言。\label{ex:1.5.2}
\begin{tnote}
原书 p.114 此处作「Let $A$ be an ideal」，据上下文应为「设 $A$ 是环」。[依据：原书 p.114 页图像，参看习题 1 的同款开头]
\end{tnote}
\begin{enumerate}[label=(\arabic*),leftmargin=2.8em]
  \item $A[t]$ 关于 $(t)$-进拓扑的完全化同构于 $A[[t]]$。
  \item $A[t]$ 关于 $(IA[t]+(t))$-进拓扑的完全化同构于 $\widehat A[[t]]$。
  \item $A[t]$ 中多项式序列 $\sum_n a_n^{(m)}t^n$（$m=0,1,\dots$）在 $(IA[t]+(t))$-进拓扑下收敛到 $0$，当且仅当对每个 $n$，$a_n^{(m)}$ 按 $A$ 上的 $I$-进拓扑满足 $\lim_{m\to\infty}a_n^{(m)}=0$。
  \item 有嵌入
  \[
  A[t]\to A^{\bN},\quad \sum_n a_n t^n\mapsto(a_0,a_1,\dots).
  \]
  考虑由 $I$-进拓扑的乘积在 $A[t]$ 上诱导的拓扑。此拓扑与 $A[t]$ 上的 $(IA[t]+(t))$-进拓扑一致。
  \item $A[t]$ 关于 $IA[t]$-进拓扑的完全化同构于
  \[
  \Bigl\{\sum_{n=0}^{\infty}a_n t^n: a_n\in\widehat A \text{ 且 } \lim_{n\to\infty}a_n=0\Bigr\},
  \]
  其中极限按 $\widehat A=\plim_n A/I^n$ 上离散拓扑空间 $A/I^n$ 的乘积拓扑取。
  \item $A[t]$ 中多项式序列 $\sum_n a_n^{(m)}t^n$（$m=0,1,\dots$）在 $IA[t]$-进拓扑下收敛到 $0$，当且仅当对每个 $n$，都有 $\lim_{m\to\infty}a_n^{(m)}=0$，且关于 $n$ 一致（按 $A$ 上的 $I$-进拓扑）。
  \begin{tnote}
  原书 p.115 此处作「if and only if for every we have $\lim_{m\to\infty}a_n^{(m)}=0$ uniformly in $n$」，其中「for every」后疑脱 $n$。[依据：原书 p.115 页图像]
  \end{tnote}
\end{enumerate}
\end{exercises}
